By 3.10 and 5.2 a), $\mathscr D$ is representable by a smooth
quasi-projective prescheme over $S$. Consider the subgroup $D$ of type
(C) of $G_{\mathscr D}$ playing the universal role relative to $G/S$;
one may then consider the functor
$\mathscr C_D:(\Sch)^\circ_{/\mathscr D}\to\Ens$ defined in terms of
the $\mathscr D$-group $D$ as $\mathscr C$ in terms of the $S$-group
$G$. One then has the following result:

\begin{proposition}{5.5}\label{XIV.5.5}
Under the preceding conditions, consider $\mathscr C$ as a functor
over the prescheme $\mathscr D$; then $\mathscr C$ is
$\mathscr D$-isomorphic to the functor $\mathscr C_D$ ``of Cartan
subgroups of $D$''.
\end{proposition}
This follows at once from the

\begin{corollary}{5.6}\label{XIV.5.6}
Let $G$ be a smooth $S$-group prescheme of finite presentation,
and $D$ a subgroup of type (C) of $G$. Then there is a one-to-one
correspondence between the Cartan subgroups of $G$ contained in
$D$ and the Cartan subgroups of $D$ (more precisely, for a
subgroup $H$ of $G$, $H$ is a Cartan subgroup of $G$ if and only
if it is a Cartan subgroup of $D$).
% typo?: kept as printed: the parenthesis does not repeat H subset D.
\end{corollary}
Indeed, this is a particular case of XII 7.9 c), taking into
account that over an algebraically closed field a subgroup
of type (C) of $G$ contains a Cartan subgroup of $G$.

For the next section, the main result obtained here is 5.2 c)
for $H$ a Cartan subgroup of $G$, which allows one to state
5.5 and thus provides a useful ``d\'evissage'' of $\mathscr C$.

\sgasection{6}{Applications to the structure of algebraic groups}
\begin{theorem}{6.1}\label{XIV.6.1}
Let $G$ be a smooth algebraic group over a field $k$. Consider
the scheme $\mathscr T$ of maximal tori of $G$, isomorphic to
the scheme $\mathscr C$ of Cartan subgroups of $G$, which is
a homogeneous space under $G^0$, and a smooth connected affine
algebraic scheme (XII 7.1 d)). Then $\mathscr C$ is a rational
variety, i.e. the field of rational functions on $\mathscr C$
is a pure extension of $k$.
\end{theorem}
\begin{proof}
We shall first give the proof in the case where $k$ is
infinite. One may clearly suppose $G$ connected, for
$\mathscr T$, hence $\mathscr C$, does not change when
$G$ is replaced by $G^0$. Moreover, by XII 7.6,
$\mathscr C$ does not change when $G$ is divided by
a central subgroup. This allows us, first dividing
by the center of $G$, to suppose $G$ affine (XII
6.1), then, dividing by its reductive center (XII
4.1 and 4.4), to suppose that the reductive center
of $G$ is trivial (XII 4.7 b)). Furthermore, we
proceed by induction on $n=\dim G$, supposing the
theorem proved for dimensions $n'<n$. If $G$ is
nilpotent, then $\mathscr C$ is reduced to a point
rational over $k$, and 6.1 is trivial. Otherwise,
the Lie algebra of $G$ is not nilpotent (1.3),
hence the Cartan subalgebras of $\mathfrak g$
have dimension $n'<n$, hence the subgroups of
type (C) of $G$ have dimension $n'<n$. Consider
then the morphism
\[
 \mathscr C\longrightarrow\mathscr D
\]
considered in 5.5. One knows by 3.10 ($k$ being
an infinite field, hence $\mathfrak g$ containing
a regular element) that $\mathscr D$ is a
rational variety, i.e. the field $K$ of rational
functions on $\mathscr D$ is a pure extension
of $k$. Consider the fiber of $\mathscr C$ at
the generic point $x$ of $\mathscr D$; by 5.5
it is the scheme of Cartan subgroups of a
certain smooth connected algebraic group $D_x$
over $K=\kappa(x)$ (namely $D_x=$ ``the generic
subgroup of type (C) of $G$''). The field $L$
of rational functions on $\mathscr C$ is
therefore isomorphic to the field of rational
functions on $\mathscr C_{D_x}$, which by the
induction hypothesis (since $\dim D_x=n'<n$)
is a pure extension of $K$. Thus by transitivity
$L$ is a pure extension of $k$.

When $k$ is finite, a different proof is needed.
One may again suppose that $G$ is affine and
connected. Note that $k$ is perfect; it follows
at once that the radical $\bar R$ of $G_{\bar k}$
is ``defined over $k$'', i.e. comes from a
subgroup $R$ of $G$. Suppose first that
$R\ne G$, i.e. $G$ is not solvable, and let
\[
 u:G\longrightarrow G'=G/R
\]
be the canonical morphism. Consider the
corresponding morphism $C\mto u(C)$
\[
 v:\mathscr C_G\longrightarrow\mathscr C_{G'}
\]
(whose definition is immediate by XII 7.1 e)).
Let $x$ be the generic point of
$\mathscr C_{G'}$; then the fiber $v^{-1}(x)$
is identified with the scheme of Cartan
subgroups of $G_K$ (where $K=\kappa(x)$)
whose image in $G'_K$ is a certain Cartan
subgroup $C'_x$ (namely, ``the generic Cartan
subgroup of $G'$''). It is therefore also
the scheme of Cartan subgroups of
$H=u_K^{-1}(C')$ (XII 7.9 c)), and since
$K$ is here an infinite extension of $k$,
it follows from the part of 7.1 already
proved that the field $L$ of rational
functions on $\mathscr C_G$, equal to
that of $\mathscr C_H$, is a purely
transcendental extension of $K$. To prove
that it is a purely transcendental extension
of $k$, it therefore suffices to prove
that this is so for $K$, i.e. one is
reduced to the case where $G$ is
semisimple. One may moreover suppose
that $G$ is adjoint (dividing $G$ by
its reductive center). But then by
3.18 one has $\mathscr C_G\simeq\mathscr D_G$,
and by 3.10 it suffices to prove that
$\mathfrak g$ admits a regular point,
which (as was noted in 3.20) is an
unpublished result of Chevalley\nde{Cf.
the Appendix, by J.-P. Serre.}.
% typo?: kept as printed: C' with no x, and reference to 7.1 rather than 6.1.
% typo: extension transcendante pur -> extension transcendante pure.

It therefore only remains to treat the
case where $k$ is finite and $G$
connected affine solvable. One in
fact has a more general result:
\begin{corollary}{6.2}\label{XIV.6.2}
Let $G$ be a smooth solvable
algebraic group over a field $k$;
then the variety $\mathscr C$ of
Cartan subgroups of $G$ is
isomorphic to an affine space
$\Spec k[t_1,\ldots,t_N]$.
\end{corollary}
One may again suppose $G$ connected
and affine. Let $G_\infty$ be the
smallest of the groups occurring
in the descending central series
of $G$ (by $G_i$ such that
$G_{i+1}=[G,G_i]$); it is thus
the smallest invariant algebraic
subgroup of $G$ such that
$G/G_\infty$ is nilpotent.
Let $C$ be a Cartan subgroup of
$G$ (one exists by 1.1); then
the image of $C$ in $G/G_\infty$
is a Cartan subgroup of it,
and is therefore equal to
$G/G_\infty$; consequently the
morphism $g\mto\Ad(g)\cdot C$
from $G_\infty$ to $\mathscr C$
is an epimorphism, and identifies
$\mathscr C$ with the homogeneous
space $G_\infty/N\cap G_\infty$,
where $N$ is the normalizer
of $C$ in $G$ (moreover equal
to $C$, as recalled in 4.4,
but that does not matter here).
Note that $U=G_\infty$ is
clearly a smooth connected
``unipotent'' algebraic group
(for over the algebraic closure
of $k$, it is contained in
the unipotent part of $G$,
by the known structure of
smooth solvable affine groups,
BIBLE 6 th. 3). When $k$ is
perfect (the only case needed
to establish 6.1), it follows
very easily that $U$ is even
$k$-unipotent, i.e. admits a
composition series by algebraic
subgroups $U_i$ such that
$U_i/U_{i+1}$ is isomorphic
to $\Ga$. In fact, Rosenlicht
has proved that this result
remains valid for a group
of the form $G_\infty$ as
above, without restriction
on $k$ (M. Rosenlicht,
\emph{Questions of rationality
for solvable algebraic groups
over non perfect fields},
Annali di Matematica 1963,
pp. 97--120, theorem 4 cor. 2),
a considerably more delicate
result which we shall admit
here. It now suffices to
apply the following lemma,
doubtless well known to
specialists:

\begin{lemma}{6.3}\label{XIV.6.3}
Let $U$ be a smooth connected
algebraic group over a field
$k$, $X=U/V$ a homogeneous
space under $U$ having a
point rational over $k$.
Suppose $U$ $k$-unipotent.
Then as a $k$-scheme, $X$
is isomorphic to an affine
space $\Spec k[t_1,\ldots,t_N]$.
\end{lemma}
Indeed, let $(U_i)_{0\leq1\leq n}$
be a composition series of
$U$ by smooth connected
subgroups, with $U_n=(e)$,
$U_0=U$, $U_i/U_{i+1}\simeq\Ga$,
the $U_i$ invariant in $U$.
Then the $K_i=U_iV$ are
algebraic subgroups of $U$
(not necessarily smooth
or connected if $V$ is
not, but that matters
little to us), and
$K_{i+1}$ is invariant
in $K_i$. One has a
canonical morphism
$U_i/U_{i+1}\to K_i/K_{i+1}$
which is an epimorphism,
which proves that
$K_i/K_{i+1}$ is either
reduced to the unit group
or isomorphic to
$\Ga/H_i$, where $H_i$
is a finite subgroup of
$\Ga$, which by Rosenlicht
(\loccit th. 2) implies
that $K_i/K_{i+1}\simeq\Ga$
(a result moreover immediate
if $k$ is perfect).
Set now $X_i=U/K_i$;
let us prove by induction
on $i$ that $X_i$ is
isomorphic to an affine
space. Indeed, if this
holds for $X_i$, let
us prove that the same
holds for $X_{i+1}$.
If $K_i/K_{i+1}=e$,
one has $X_i=X_{i+1}$,
and this is trivial.
Otherwise, $X_{i+1}$ is
a principal bundle with
base $X_i$ and structural
group
$\Ga\simeq K_i/K_{i+1}$.
Thus $X_i$ being affine,
$X_{i+1}$ is a trivial
bundle, hence isomorphic
to $X_i\times\Ga$,
which again proves that
$X_{i+1}$ is isomorphic
to a standard affine
space. This proves 6.2
and consequently finishes
proving 6.1.
% typo?: kept as printed: 0 <= 1 <= n in the subscript.
\end{proof}

\begin{corollary}{6.4}\label{XIV.6.4}
Let $G$ be a smooth
algebraic group over
an infinite field $k$.
Then the set of points
of $\mathscr C$ (notation
of 6.1) rational over
$k$ is dense for the
Zariski topology.
The union of the
Cartan subgroups of
$G$ is dense in $G$.
\end{corollary}
The first assertion
holds for every
unirational variety
over an infinite
field, and is moreover
here for us the most
important ``arithmetic''
consequence of the
unirationality results.
The second assertion
follows from the
first and from the
density theorem XIII 2.1.

\begin{corollary}{6.5}\label{XIV.6.5}
Let $G$ be a smooth
connected algebraic
group over $k$.
Then the variety $Z$
of regular semisimple
points of $G$ (XIII
3.5) is a unirational
variety. In particular,
if $k$ is infinite,
the set of points
of $Z$ rational over
$k$ is dense in $Z$.
\end{corollary}
Indeed, $Z$ is an
open subset of a
torus over
$\mathscr C$, hence
its function field
$L$ is the function
field of a torus
defined over the
function field $K$
of $\mathscr C$
(namely the ``generic
maximal torus'' of
$G$); it is therefore
a unirational extension
of $K$ (XIII 3.4),
and since $K$ is a
pure extension of
$k$ by 6.1, $L$
is a unirational
extension of $k$.

\begin{corollary}{6.6}\label{XIV.6.6}
Let $G$ be a smooth
connected algebraic
group over $k$,
and let $H$ be
the subgroup of
$G$ generated by
the subvariety $Z$
of regular semisimple
points, i.e. (XII
8.2) the smallest
invariant algebraic
subgroup of $G$
such that $G/H$
has reductive rank
zero (i.e. over
the algebraic closure
of $k$ is an
extension of an
abelian variety by
a smooth connected
unipotent group).
Then $H$ is a
unirational variety.
In particular, if
$G=H$, i.e. (XII
8.4) if $G$ is
affine and if over
the algebraic closure
$\bar k$ there
exists no nontrivial
homomorphism from
$G_{\bar k}$ to
$\Ga$, then $G$
is a unirational
variety, hence if
$k$ is infinite,
the set of its
points rational over
$k$ is dense.
\end{corollary}
This follows at
once from 6.5,
for it is immediate
that if one has
smooth connected
$k$-preschemes $Z_i$
which are unirational
varieties and
morphisms
$u_i:Z_i\to G$,
then the algebraic
subgroup of $G$
generated by the
$u_i$ (VI\textsubscript{B}
7.1) is a
unirational variety.

As an interesting
particular case of
6.5 or 6.6 (as
one prefers), note:
\begin{corollary}{6.7}\label{XIV.6.7}
Let $G$ be a smooth
connected affine
algebraic group of
unipotent rank zero;
then $G$ is a
unirational variety.
\end{corollary}
One may make 6.4
more precise as
follows:
\begin{corollary}{6.8}\label{XIV.6.8}
Let $G$ be a smooth
connected algebraic
group over the
infinite field $k$,
and $C$ a Cartan
subgroup of $G$.
Then the union of
the conjugates of
$C$ by regular
semisimple elements
of $G(k)$ is dense
in $G$.
\end{corollary}
This follows at
once from 6.4
and XIII 3.6,
which says that
the morphism
$\varphi:Z\times C\to G$
defined by
$\varphi(t,c)=\ad(t)c$
is dominant.
This result also
implies (without
supposing $k$
infinite):
\begin{corollary}{6.9}\label{XIV.6.9}
Let $G$ be a smooth
connected algebraic
group over $k$,
$H$ a smooth
connected algebraic
subgroup of $G$
such that $H$ has
the same reductive
rank and the same
nilpotent rank as
$G$ (i.e. over
the algebraic closure
$\bar k$ of $k$,
$H_{\bar k}$ contains
a Cartan subgroup
of $G_{\bar k}$).
If $H$ is a
unirational variety,
so is $G$.
If $H(k)$ is
dense in $H$,
$G(k)$ is dense
in $G$.
\end{corollary}
Indeed, the morphism
$\varphi:Z\times H\to G$
defined by
$\varphi(t,h)=\ad(t)h$
is dominant; now
by 6.5, $Z$ is
a unirational variety,
and by hypothesis
so is $H$, hence
$Z\times H$, whence
the first result.
The second is
proved similarly.

We now recover
the following
well-known result
(due to Chevalley
in characteristic
$0$, to Rosenlicht
in characteristic
$p>0$):
\begin{corollary}{6.10}\label{XIV.6.10}
Let $G$ be a smooth
connected affine
algebraic group over
a perfect field $k$.
Then $G$ is a
unirational variety,
hence if $k$ is
moreover infinite,
$G(k)$ is dense
in $G$.
\end{corollary}
Indeed, by 1.1,
$G$ admits a Cartan
subgroup $C$. By
6.9 it suffices
to prove that the
latter is a
unirational variety.
Now $k$ being
perfect, one sees
at once by Galois
descent from the
case of algebraically
closed $k$ (BIBLE
6 th. 2) that
one has $C=T\times C_u$,
where $T$ is the
maximal torus of
$C$ and $C_u$
a smooth connected
unipotent group.
One already knows
that $T$ is a
unirational variety
(XIII 3.4); it
remains to see
that the same
holds for $C_u$.
Now $k$ being
perfect, $C_u$ is
even $k$-unipotent,
and one may
apply 6.3.

\begin{remarks}{6.11}\label{XIV.6.11}
\noindent a) There
are known examples
(Rosenlicht) of
twisted forms of
$\Ga$ over an
imperfect field
which have only
finitely many
rational points,
hence a fortiori
are not unirational
varieties. On the
other hand, Chevalley
has given an
example of a torus
over a field of
characteristic zero
which is not a
rational variety.
On the other hand,
it follows from
Chevalley's theory
of semisimple groups
that over an
algebraically closed
field, every smooth
connected affine
algebraic group is
a rational variety.
Note moreover that
the question of
unirationality arises
in any event only
for affine algebraic
groups, a unirational
algebraic group being
necessarily affine
by Chevalley's
structure theorem.

\noindent b) With
the notation of
6.6, it is tempting
to try to give
a condition of
unirationality of
$G$ in terms of
the group $G/H$
(which is unipotent
if $G$ is assumed
affine). The latter
must clearly be
unirational; is
this condition also
sufficient? Note
that an example
of Rosenlicht
(\loccit) shows
that a smooth
connected unipotent
algebraic group $U$
can be a rational
variety without
being $k$-unipotent.

\noindent c) It
would be interesting
to study, over a
finite field $k$,
questions of the
``density'' type,
such as the
following (raised
by Rosenlicht):
Let $G$ be a
smooth connected
algebraic group
over $k$; is $G$
then generated by
its Cartan subgroups%
{\renewcommand{\thefootnote}{*}\footnote{This
question has since
been answered
in the affirmative
by Steinberg.}\addtocounter{footnote}{-1}}%
\nde{Having been unable to identify this result in the
\emph{Collected Papers} of R. Steinberg, let us give a proof based
on the Bruhat decomposition. Let $G/k$ be a semisimple group
defined over the finite field $k$. One must show that $G$ is
generated by its $k$-Cartan subgroups. This question is stable
under central isogeny and restriction of scalars; one is
therefore reduced to the case where $G$ is geometrically
simple (in the same way as in lemma 1 of the appendix below).
Denote by $G^\sharp$ the subgroup of $G$ generated by its
$k$-subtori. The group $G$ is quasi-split (Lang) and therefore
admits a Killing pair $(T,B)$. By definition, one has
$T\subset G^\sharp$; in particular $T$ normalizes $G^\sharp$.
The big cell $R_u(B^-)\,T\,R_u(B)$ of $G$ indicates that it
suffices to verify that $R_u(B)\subset G^\sharp$. Let $S$ be
the maximal split $k$-torus of $T$. One considers the relative
root system $\Phi(G,S)$ and a base $\Delta_k$. Given
$\alpha\in\Phi(G,S)$, one denotes by $U_{(\alpha)}$ the unipotent
subgroup associated to $\alpha$ (cf. A. Borel, \emph{Linear
Algebraic Groups}, second edition (1991), Springer, Prop.
21.9). Since the $k$-group $R_u(B)$ is generated by the
unipotent $k$-groups $U_{(\alpha)}$ ($\alpha\in\Delta_k$), one
is reduced to verifying that $U_{(\alpha)}\subset G^\sharp$.
A glance at the classification shows that there exists a
semisimple group $G_\alpha$ of quasi-split type $A_1$, ${}^2A_1$,
${}^3A_1$ or ${}^2A_2$ such that $U_{(\alpha)}\subset G_\alpha\subset G$.
It is therefore permissible to suppose that $G=\PGL_2$ or
$G=\operatorname{SU}_3(K)$, where $K$ denotes the unique quadratic
field extension of $k$. The group $\PGL_2^\sharp$ containing the
standard split torus $T$, the possibilities up to conjugacy
under $G(k)$ are the following: $\PGL_2^\sharp=T$,
$\PGL_2^\sharp=B$, or $\PGL_2^\sharp=\PGL_2$. The case
$\PGL_2^\sharp=T$ is excluded since $\PGL_2^\sharp$ contains
the $k$-torus $R_{K/k}(\Gm)/\Gm$. The preceding discussion
indicates that if $\PGL_2^\sharp=B$, then $\PGL_2^\sharp=\PGL_2$.
One therefore has $\PGL_2^\sharp=\PGL_2$. If
$G=\operatorname{SU}_3(K)$, $\Delta_k=\{\alpha\}$ and the
possibilities for $G^\sharp$ (up to conjugacy) are the
following: $G^\sharp=T$, $G^\sharp=B$,
$G^\sharp=U_{(2\alpha)}\rtimes T$,
$G^\sharp=\langle U_{(2\alpha)},U_{(-2\alpha)},T\rangle=\SL_2\cdot T$,
or $G^\sharp=G$. The case $G^\sharp=B$ is excluded as for
$\PGL_2$. The case $\SL_2\cdot T$ is excluded because
$G^\sharp$ contains the $k$-torus
$R^1_{K/k}\bigl(R^1_{K_3/K}(\Gm)\bigr)$, $K_3/K$ denoting the
extension of degree $3$ of $K$, and the isogeny class of
this torus is irreducible. Moreover, this torus also
excludes the cases $G^\sharp=T$ and
$G^\sharp=U_{(2\alpha)}\rtimes T$. One concludes that $G^\sharp=G$.}?

In this question, one can reduce to the case of affine $G$,
by dividing by the center. The answer would be affirmative
in the semisimple case if one could make the result of
Chevalley on existence of regular points mentioned in 3.20
more precise, so as to obtain a regular element of
$\mathfrak g$ which does not belong to $\mathfrak h=\Lie(H)$,
where $H$ is a smooth algebraic subgroup of $G$ with
$H\ne G$ given in advance ($G$ an adjoint semisimple group
over the finite field $k$).
\end{remarks}
